% Início do documento, define parâmetros principais a serem seguidos.
\documentclass[
    % opções da classe memoir
    article,
    12pt,
    a4paper,
    % opções da classe babel
    english,
    brazil,
    sumario=tradicional
]{/home/lucas_galdino/Desktop/TCC/planning/abntex2/abntex2}

% Pacotes básicos

\usepackage{lmodern}			% Usa a fonte Latin Modern
\usepackage[T1]{fontenc}		% Selecao de codigos de fonte.
\usepackage[utf8]{inputenc}		% Codificacao do documento (conversão automática dos acentos)
\usepackage{indentfirst}		% Indenta o primeiro parágrafo de cada seção.
\usepackage{nomencl} 			% Lista de simbolos
\usepackage{color}				% Controle das cores
\usepackage{graphicx}			% Inclusão de gráficos
\usepackage{microtype} 			% para melhorias de justificação


% Pacotes de citação

\usepackage[brazilian,hyperpageref]{backref}	 
% Paginas com as citações na bibl. "brazilian" siginfica que seguirá os padrões brasileiros de citação, como datas, caracteres especiais, etc.
% hyperpageref deixa o número da página como hiperlink
\usepackage[alf]{/home/lucas_galdino/Desktop/TCC/planning/abntex2/abntex2cite}	
% Citações padrão ABNT

% my own changes to the code
\usepackage{/home/lucas_galdino/Desktop/TCC/planning/abntex2/my_changes}
%


% Informações de dados para CAPA e FOLHA DE ROSTO
\titulo{Análise cŕitica do TCC Desenvolvimento de Método Computacional Para Definição de Rotas Baseadas No Uso de Sistemas de Transporte Coletivo}
\autor{
    Lucas Eduardo Galdino Silva\thanks{Universidade Federal de Santa Catarina.\ E-mail: lucasegaldinos@gmail.com}
    }
\local{Brasil}
\data{2024}

% Setando asa informações do PDF
\makeatletter % makeatletter allows to use @ in the command names
\hypersetup{ % hypersetup requires this, as it is used to set the PDF metadata and needs to have the properties fully expanded. This does not happen if the command is used directly as `\titulo`.
    pdftitle={\@title}, % Expands the title macros
    pdfauthor={\@author}, % Expands the author macros
    pdfsubject={Análise cŕitica do TCC Desenvolvimento de Método Computacional Para Definição de Rotas Baseadas No Uso de Sistemas de Transporte Coletivo},
    pdfkeywords={Análise cŕitica, TCC, Desenvolvimento de Método Computacional, Definição de Rotas, Sistemas de Transporte Coletivo},
    pdfcreator={\@author}, % Expands the author macros,
    bookmarksdepth=4 % The depth of the bookmarks in the PDF
}
\makeatother % makeatother restores the normal meaning of @




% ---
% Document start

\begin{document}

% Set language
\selectlanguage{brazil}
% Remove obsolete space from phrases
\frenchspacing

% Set title
\maketitle

\begin{resumoumacoluna}
    O documento realizará uma análise crítica do Trabalho de Conclusão de Curso (TCC)
    intitulado \cite{tcc_vitorhp}.
    O objetivo é avaliar a qualidade do trabalho, identificar suas contribuições, limitações e oportunidades para futuros estudos.
    A análise será conduzida com base em critérios de avaliação acadêmica, considerando a clareza, coerência, relevância e originalidade do trabalho.
    O documento apresentará uma visão geral do TCC, uma análise detalhada das partes constituintes, como introdução, objetivos, metodologia, resultados
    e discussão, e uma conclusão com as principais contribuições e sugestões para trabalhos futuros.

    \vspace{\onelineskip}

    \noindent
    \textbf{Palavras-chave}: latex.\ abntex.\ editoração de texto.
\end{resumoumacoluna}

\textual

\section{Introdução}

O TCC foi escolhido devido ao meu interesse e familiaridade em relação ao tema
otimização de rotas. Também tenho interesse em desenvolver algo parecido, mas abordando outro problema relacionado
ao famoso Problema do Caixeiro Viajante.


\section{Análise das Partes Constituintes}

Serão analisados aqui diversos aspectos da abordagem do estudante Victor
Hammann Pereira.

\subsection{Partes Constituintes}

Ao iniciar o trabalho o autor apresenta
\begin{enumerate}
    \item \textbf{Resumo}: O trabalho possui um resumo, explicitando o que será discutido dentro do trabalho. O resumo é claro e objetivo,
          apresentando uma visão geral do problema, da metodologia utilizada e dos resultados obtidos.
    \item \textbf{Abstract}: O texto possui também um abstract, que é uma versão resumida do trabalho em inglês. O abstract é claro e objetivo, apresentando uma visão geral do problema, da metodologia utilizada e dos resultados obtidos.
    \item \textbf{Lista de siglas}: As siglas utilizadas no trabalho.
    \item \textbf{Lista de figuras}: Uma lista com todas as figuras utilizadas no trabalho, com seus respectivos títulos e páginas.
    \item \textbf{Lista de tabelas}: Uma lista com todas as tabelas utilizadas no trabalho, com suas legendas e páginas.
    \item \textbf{Sumário}: Contendo as informações sobre os capítulos do trabalho.
\end{enumerate}\par

Em seguida, no primeiro capítulo, são apresentados
\begin{enumerate}
    \item \textbf{Justificativa}: Logo na introdução, o autor apresenta a justificativa do trabalho, explicando a importância do problema e a necessidade de uma solução computacional para a definição de rotas baseadas no uso do sistema de transporte coletivo.
    \item \textbf{Objetivos}: O autor define os objetivos gerais e os objetivos específicos do trabalho.
    \item \textbf{Considerações iniciais}: O autor apresenta as considerações iniciais, explicando as limitações da metodologia adotada.
\end{enumerate} \par

No segundo capítulo, são apresentados
\begin{enumerate}
    \item \textbf{Revisão bibliográfica}: São analisados trabalhos similares ou dentro do proposto pelo autor
    \item \textbf{Materiais}: São apresentados os materiais utilizados para o desenvolvimento do trabalho, como linguagem de programação utilizada, serviço de bancos de dados, bibliotecas e ferramentas GIS (Geographic Information System).
\end{enumerate}\par

No terceiro capítulo, o autor descreve a \textbf{metodologia} utilizada para o desenvolvimento do trabalho, iniciando através da descrição de informações necessárias, preparação dos dados e desenvolvimento do modelo. \par

O quarto capítulo, são apresentados os \textbf{resultados} obtidos pelo autor através da metodologia proposta. \par

No quinto capítulo, são apresentadas \textbf{conclusões} do trabalho, resumindo os principais aprendizados, resultadoes e implicações.
São desenvolvidas também recomendações para futuros trabalhos. \par

Finalmente o autor apresenta
\begin{enumerate}
    \item \textbf{Referências bibliográficas}: O trabalho possui uma lista das referências bibliográficas citadas ao longo do texto.
    \item \textbf{Apêndices}: Nos apêndices são apresentados os códigos desenvolvidos pelo autor.
\end{enumerate}\par

\section{Análise Elaborada das Partes}

\subsection{Análise da Justificativa}

O autor inicia a justificativa explicando como o transporte permite às pessoas
se deslocarem sem a necessidade de possuir um automóvel. Em seguida a presenta
os problemas enfrentados para a adoção do mesmo, sendo a apresentação do mesmo
um dos principais problemas a ser resolvido. Uma solução para este problema aumentaria a
adoção e a qualidade dos serviços ofertados. Ainda cita problemas relacionados ao uso
excessivo de carros.\par

A justificativa expõe alguns problemas relacionados ao serviços de transporte coletivo.
O autor, no entanto poderia trazer outros problemas em que a solução proposta pode ser aplicada,
como por exemplo, a otimização do tempo de viagem dos usuários, a redução de emissões de gases de
efeito estufa, a melhoria da segurança dos usuários, entre outros.\par

\subsection{Análise dos Objetivos}

Os objetivos gerais, sintetizam o propósito do trabalho quanto a utilização
de softwares abertos para o roteamento de ônibus utilizando softwares abertos,
porém peca ao desconsiderar o objeto de estudo como a cidade de Florianópolis ao generalizar para a
determinação em outros locais. No entanto, cita tais problemas ao falar das limitações do modelo.
Tal solução só seria possível levando em consideração as diversas características e particularidades de cada local.\par

Alguns dos objetivos específicos são bem definidos e claros. O autor objetiva identificar as ferramentas disponíveis para
fornecimento de informações relativos a rotas, tais como softwares GIS e aplicar funcionalidades destes para processamento
de dados espaciais relativos a infraestrutura do transporte. Aplicar também técnicas de criação e modelagem de bancos de dados,
coleta e limpeza de dados para o suporte do serviço e analisar os resultados obtidos.\par

O autor ainda cita as limitações do seu modelo, como a área de estudo sendo limitada a Florianópolis, a desconsideração de fatores
externos ao sistema de transporte coletivo, como não haver o acesso em tempo real aos ônibus, a desconsideração de trânsito e curvas
— o cálculo é utilizado calculando pontos absolutos de distância —, o cálculo do tempo ser baseado em estimativas fornecidas pela
própria prestadora de serviços (estas últimas, ao meu ver, as que mais prejudicam a aplicação), a desconsideração de fatores de custos
de passagens aos usuários do sistema de transporte.\par

Os objetivos não explicitam que tipo de método ou qual problema dentro do VRP busca ser resolvido, quais técnicas pretende implementar,
e difere da justificativa do problema. O objetivo está contextualizado com ênfase em ferramentas e não no problema em si.\par

\subsection{Análise da Conclusão}

O trabalho desenvolvido satisfez o objetivo do autor em identificar as ferramentas disponíveis para fornecimento de informações
relativos a rotas, utilizando diferentes softwares abertos de GIS, bancos de dados e para cálculos computacionais.
Os principais softwares abertos utilizados foram o:
\begin{enumerate}
    \item OSM (OpenStreetMap) para a coleta de dados de infraestrutura de transporte;
    \item QGIS e GRASS como softwares de GIS para a visualização e análise de dados espaciais;
    \item PostgreSQL e PostGIS com a extensão pgRouting para a criação e modelagem de bancos de dados espaciais;
    \item Python para o desenvolvimento do algoritmo de cálculo de rotas.
\end{enumerate}\par

O autor realiza o objetivo de utilizar softwares abertos para o desenvolvimento, destacando os principais problemas encontrados durante o processo e ressaltando que diversas simplificações tiveram que ser realizadas, principalmente quando se tratando do cálculo de rotas. Destaca também que não foram levadas em consideração o sentido das linhas de ônibus, piorando os resultados obtidos, e que, melhores resultados poderiam ser obtidos, mas a eliminação de resultados ruins seriam uma tarefa mais complexa.
Apesar destes problemas, a saída através da linha de comando é bem informativa, mesmo não permitindo a visualização através do mapa.\par

No geral, o objetivo principal proposto pelo trabalho foi atingida, mesmo que o discente não tenha obtido um produto realmente funcional.\par

\subsection{Análise do Corpo do Trabalho}

O problema a ser resolvido é uma variação do famoso Problema do Caixeiro Viajante (TSP - Traveling Salesman Problem), que consiste em encontrar o caminho mais curto possível que visita um conjunto de cidades e retorna à cidade de origem. O nome da variação, apesar de não citado, é denominada Problema de Roteamento de Veículos Capacitado com Coleta, Entrega e Janelas de Tempo (CVRPPDTW - Capacitated Vehicle Routing Problem with Pickup, Delivery and Time Windows), que adiciona restrições de capacidade dos veículos e janelas de tempo de atendimento aos clientes.\newline
Este problema é um problema NP-difícil, o que significa que não existe uma solução eficiente conhecida para resolvê-lo em tempo polinomial. Portanto, é necessário utilizar técnicas de otimização, como algoritmos de busca local, algoritmos genéticos ou métodos heurísticos para encontrar soluções aproximadas. É largamente estudado, principalmente na área de pesquisa operacional, otimização e IA. Tem aplicações em diversas áreas, como logística, transporte, engenharias, eletrônica e ciência da computação. Estas observações não são trabalhados no texto, mas seriam importantes para contextualizar o problema e mostrar a relevância da solução proposta.\par

A justificativa do trabalho apresenta a importância do problema, mas não entra nos detalhes sobre necessidade de uma solução computacional para a definição de rotas baseadas no uso do sistema de transporte coletivo para o usuário do sistema.\par
O autor poderia ter citado a relevância do
problema na justificativa, explicando como o uso do sistema de transporte coletivo pode contribuir para a redução de emissões de gases de efeito estufa e reduzir o tempo de deslocamento.\par

A busca por uma solução computacional para lidar com a complexidade dos dados envolvidos, como a grande quantidade de informações sobre os sistemas de transporte coletivo, a demanda de passageiros e as restrições de capacidade dos veículos. Tal solução, se bem aplicada, também diminuiria o tempo de planejamento, necessidade de recursos humanos para a definição de rotas (consequentemente diminuindo erros humanos), melhoraria a satisfação dos passageiros, diminuiria o uso de recursos na mesma linha, teoricamente diminuiria gastos públicos com empresas terceirizadas e manutenção de vias e ainda contribuiria para a redução de congestionamentos nas cidades. Essas são apenas algumas das possíveis aplicações do trabalho, que podem ser ainda mais impactantes e que falhei em observar.\par

O trabalho realizado possui todas as partes constituintes necessárias, onde o que não consta, não era aplicável. A metodologia utilizada é bem explicada, a ênfase dada ao uso de softwares abertos é positiva, e o discente demonstra um bom conhecimento adquirido através do processo apesar das simplificações feitas. Julgo tais simplificações normais, pois tal desenvolvimento iria além do escopo proposto por um trabalho de conclusão de curso, necessitando de ferramentas mais robustas, métodos não tão acessíveis, parcerias com empresas privadas para a obtenção de dados precisos, financiamento e tempo. A revisão bibliográfica também é bem trabalhada, trazendo diversas referências e resultados de trabalhos já desenvolvidos.\par

\subsection{Análise dos resultados}

Devido às simplificações e complexidade do tema proposto, muitas simplificações são realizadas, gerando resultados que apesar de realizarem o proposto, não são tão satisfatórios.\par

\section{Análise da Língua Portuguesa}

A partir do que foi lido, não identifiquei nenhum erro ortográfico crasso ou não conformação às normas da língua. O texto é bem escrito e o discente demonstra um bom conhecimento da língua.

\section{Comentários Finais}
O trabalho apresenta limitações, mas estas são normais para um trabalho de conclusão de curso
e não prejudicam o objetivo principal proposto.\newline
Ao meu ver, trata-se de um ótimo trabalho, com boa revisão bibliográfica. O discente demonstra
ter adquirido bastante conhecimento dos softwares utilizados.\newline
Além disso, o uso de softwares abertos democratiza a acessibilidade a quem possa se interessar pelo tema.\par

O autor também disponibiliza o código fonte do trabalho,
o que é um ponto positivo, pois permite que outros possam replicar o resultado,
e até mesmo melhorá-lo. Seria interessante se o banco de dados
utilizado estivesse disponibilizado em anexo, mas é factível que isso não seja possível
devido ao seu tamanho, tendo o discente que arcar com os custos da disponibilidade, tendo o mesmo provido as
informações necessárias para que a replicação do trabalho seja possível.\par

\bibliography{refs.bib}

\end{document}